% Proposed revision of the whole Methods section of the Overleaf main.tex (main.tex lines 157-208 as of 29 October
% 2026), checked against the Snow-polygon primary v13 (primary_snow520_merge_dhsmics_gamma2_v13; sensitivities v7;
% imputed-covariate version primary_snow520_merge_imputed_gamma2_v12). Built on the author's current text; supersedes
% methods_data_sources_proposed.tex and methods_covariate_adjustment_proposed.tex. main.tex is not edited; paste by hand.
% New citation keys needed in ref.bib: igme (UN IGME all-cause estimates), wpp (UN World Population Prospects).
\section*{Methods}
\subsection*{Data sources}

Complete birth histories were obtained from the Births Recodes of the Demographic and Health Surveys and Malaria Indicator Surveys (DHS/MIS) \cite{dhs} and from UNICEF Multiple Indicator Cluster Surveys (MICS) that recorded a complete birth history \cite{mics}. Birth dates, survival status at the time of the survey interview, and reported age at death were used to construct individual child trajectories. Ages were grouped in 7 age-bands: $<1$ month, 1--5 months, 6--11 months, 1--2 years, 2--3 years, 3--4 years and 4--5 years. An age band was eligible to be included in the analysis if the child entered it alive in the 60 months before the interview and in a calendar year between 2000 and 2024, and the whole band ended before the interview, irrespective of the survival outcome.

The units of analysis were the 520 malaria-endemic sub-national administrative polygons of sub-Saharan Africa defined by Snow et al.\ \cite{snow:17}. Each survey record was located by the lowest-level administrative unit reported in its survey: the survey region or, where a survey region spanned several polygons, a finer location variable (for example state, province or district) matched by boundary geometry, polygon name or a reviewed alias. A record was assigned to the polygon holding at least 90\% of the population of that location (Gridded Population of the World version 4 (GPWv4), 2020 \cite{gpwv4}) or, failing that, to the union of the polygons each holding at least 5\%; locations with more than half of their population in malaria-free or island polygons were excluded. We refer to a survey's records in one such unit as a survey unit; the analysis used 1,452 survey units covering 465 polygons, of which 215 were unions of two or more polygons.

Thirteen area-level covariates were computed for each survey unit from the survey microdata, using the DHS indicator definitions and the survey sampling weights. From the DHS/MIS Births Recodes: the mean maternal age at first live birth, mean maternal education in completed years, mean household wealth quintile and percentage of urban residence among mothers with a live birth in the 60 months before interview; coverage of three doses of a DTP-containing vaccine (DTP3) and of a measles-containing vaccine among living children aged 12--23 months; the percentage of births in the 60 months before interview delivered in a health facility; and the percentage of non-first births with a preceding birth interval shorter than 24 months. From the household and height-and-weight recodes: the percentage of households with an improved drinking-water source, an improved sanitation facility and electricity, and the prevalence of wasting and stunting (below $-2$ SD of the WHO Child Growth Standards) among de facto children under five years. National values computed in the same way were checked against the published DHS estimates. Missing DTP3 or measles coverage was replaced by the WHO/UNICEF national estimate (WUENIC, 2025 revision) for the survey year, and any covariate still missing in a survey unit took the mean of the survey's other units, weighted by their number of records.

National child HIV incidence estimates for ages 0--14 years, expressed as new infections per 1,000 uninfected population, were obtained from the UNAIDS 2025 estimates distributed by UNICEF. For Nigeria, which has no child series, child incidence was imputed from a model informed by the corresponding adolescent (15--19 years) incidence series (see Covariate adjustment); Liberia, which has no UNAIDS/UNICEF incidence series, was assigned the rate implied by UNAIDS counts of new infections in children aged 0--14 years. National GDP per capita, current health expenditure per capita (both in current US\$) and political stability and absence of violence were obtained from the World Bank World Development Indicators and Worldwide Governance Indicators \cite{wdi,wgi}. These four national covariates were matched to the country and calendar year of band entry. For years before South Sudan's independence in 2011, South Sudan took Sudan's values for GDP, health expenditure and political stability; and because the Worldwide Governance Indicators were published every two years until 2002, political stability in 2001 was the mean of each country's 2000 and 2002 values. Health expenditure for 2024, not yet published for any country in the analysis, was the country's 2023 value.

Of the 93 MICS surveys from sub-Saharan Africa, 47 (28 countries) recorded a complete birth history; the other 46 recorded only summary birth histories, which cannot be used in the age-band analysis, and were excluded. Each MICS birth history was converted to the DHS Births Recode layout and processed with the same age-band construction and eligibility rules as the DHS/MIS data. MICS does not publish region boundaries, so the 366 MICS region labels were matched, by exact name (291) or a reviewed correspondence table (75), to the regions of a DHS survey in the same country; where no DHS boundaries exist, they were matched to admin-1 polygons (Guinea-Bissau and the Central African Republic) or directly to the Snow et al.\ polygons (South Sudan). The 13 covariates were computed from the MICS files with the same definitions; education, recorded as level and grade, was converted to completed years, and facility delivery refers to the last birth in the two years before interview.

\emph{P.\ falciparum} prevalence estimates were taken from the Malaria Atlas Project (MAP) annual modelled surfaces of the parasite rate standardised to ages 2 to 10 years (\textit{Pf}PR$_{2-10}$) for 2000--2024 \cite{map,map2025}. Each polygon's yearly estimate was the mean of the surface over the polygon, weighted by GPWv4 2020 population counts (population density $\times$ cell area) on the MAP grid; the value for a union of polygons was the population-weighted mean of its polygons \cite{gpwv4}. Each age band was assigned the \textit{Pf}PR$_{2-10}$ of its unit in the calendar year of entry into the band. For national burden calculations, country-level prevalence was calculated using population-weighted means, with weights derived from the GPWv4 population density surface for 2020 and grid-cell areas \cite{gpwv4}.

For the burden calculations, national all-cause deaths at $<1$ month, 1--11 months and 1--4 years were obtained from the UN Inter-agency Group for Child Mortality Estimation (UN IGME) \cite{igme}, with the UN World Population Prospects population under 5 as the child-years at risk \cite{wpp}. A sensitivity analysis used IHME Global Burden of Disease all-cause deaths and person-years by age \cite{ihme}. Comparator estimates of under-five malaria deaths were obtained from the CA-CODE 2026 series disseminated through the UN IGME portal \cite{Perin:26,Villavicencio:24} and from IHME.


\subsection*{Covariate adjustment}

Covariate selection was guided by the simplified causal processes shown in Figure \ref{fig:dag}. This causal model was used to guide the selection of variables that are common causes of transmission intensity and mortality and that lie off the causal pathway. These are primarily variables that are proxies for socio-economic status and general health-system quality. We also selected variables that are determinants of childhood mortality from other causes (not necessarily confounders, but variables that explain variation in all-cause mortality). These include vaccination coverage, child HIV incidence and delivery in a health facility. We did not adjust for long-lasting insecticide treated net (LLIN) coverage (measured in DHS/MIS surveys) because this is a determinant of prevalence (LLINs reduce transmission) and corresponding reductions in mortality are mediated by the reduction in transmission. For malnutrition, causality could go in both directions and so measures of malnutrition could mediate the relationship between \textit{Pf}PR$_{2-10}$ and death, or confound it. To assess the potential impact of this, we carried out a sensitivity analysis with and without measures of malnutrition (wasting and stunting).

Adjustment included 17 covariates: the 13 area-level summaries and 4 national measures described above. HIV incidence, GDP and health expenditure were natural-log transformed. All covariates were centred and scaled by their mean and standard deviation over the analysis sample, and entered linearly, with separate coefficients in each age-band model. Wealth quintiles are defined within each survey, so the mean wealth quintile measures relative position within the country. The area-level summaries refer to the survey (household conditions and anthropometry at interview; births, deliveries and vaccinations in the preceding one to five years) and take the same value for every age band of that survey, whereas the national covariates vary with the year of band entry.

Summaries were calculated from all respondents in each survey unit, before complete-case selection. Births Recode summaries used the women's sampling weights, household and anthropometric summaries the household weights, and MICS summaries the MICS women's, household and children's weights. Questions not asked in a survey were excluded from the denominator; following the DHS tabulation rules, missing or `don't know' answers on vaccination and place of delivery were counted as not vaccinated and not delivered in a facility, and other missing responses were excluded.

Child HIV incidence for Nigeria was imputed with a Bayesian model of log incidence informed by adolescent incidence and temporal and regional patterns; the primary analysis used the posterior medians. Records with remaining missing covariates were excluded (822,434 records and 13,296 deaths), mainly because of survey-wide gaps in wasting and stunting, facility delivery, wealth or electricity, and missing health expenditure for Somalia and for Zimbabwe before 2010. Imputation uncertainty was not propagated in the primary analysis (see Sensitivity analyses).


\subsection*{Statistical model and model selection}

The unit of observation was the child age-band record, with the eligibility rules above \cite{Wakefield:19,Burstein:18}. The outcome was death versus survival during the age band. Because a band was included only if it could have been completed by the interview date, a death in a band cut short by the interview was not counted (for example, a child who died a month before the interview at age 4.5 years contributed only to the earlier age bands). A child surviving to five years therefore contributed up to seven records; a child dying at eight months contributed three. Because the exposure and all covariates are defined at the level of the survey unit and calendar time, the likelihood of the individual records is equivalent to a binomial likelihood for the number of deaths among children entering each age band in each survey unit and entry month.

We fitted seven separate models, one per age band, because both the effect of \textit{Pf}PR$_{2-10}$ and the effects of the confounding variables could differ by age band. Each model was a binomial generalised additive model with a complementary log-log link and an offset equal to the log band width in years, so that the linear predictor is the log annualised mortality hazard, giving a proportional-hazards model for grouped survival times.

The linear predictor comprised a penalised cubic regression spline of the unit's \textit{Pf}PR$_{2-10}$ at band entry (basis dimension 5), a penalised cubic regression spline of calendar time at band entry (basis dimension 6), the 17 standardised covariates entered linearly, and independent Gaussian random intercepts for survey, country and survey unit. Spline knots were fixed in advance: five \textit{Pf}PR$_{2-10}$ knots spanning the observed range (0--95\%) and six calendar-time knots evenly spaced from January 2000 to December 2023 (records entering a band in 2024, 0.1\% of the sample, lie beyond the last knot, where the spline is linear). Models were fitted by fast REML with discretised covariates using \texttt{bam} in the R package \texttt{mgcv} (version 1.9-4), with $\gamma = 2$ to give smoother curves. Hazard ratios between two prevalence levels $X$ and $Y$ were calculated as $\exp\{f_g(Y) - f_g(X)\}$ from the fitted spline $f_g$ of age band $g$. Confidence intervals used the coefficient covariance conditional on the estimated smoothing parameters and the fixed HIV imputation. Full model equations are given in the Supplement (SText~1).

\subsection*{Malaria-attributable mortality}

For each country, year and age band $g$, the malaria-attributable fraction of all-cause mortality was calculated as $1 - \exp\{f_g(0) - f_g(P)\}$, where $P$ is the population-weighted national mean \textit{Pf}PR$_{2-10}$ for that year, and multiplied by that band's all-cause deaths. The fitted curves were applied to the 42 sub-Saharan African countries with national MAP estimates, including those not contributing surveys to the model, annually for 2000--2024. Malaria-attributable mortality rates were these deaths, summed over age bands, per 1,000 child-years under 5. In the primary calculation, UN IGME deaths at 1--11 months and 1--4 years were split across the model's age bands in proportion to the deaths observed in the survey sample (1--5 and 6--11 months: 52\% and 48\%; 12--23, 24--35, 36--47 and 48--59 months: 38\%, 32\%, 19\% and 11\%), and child-years were the UN WPP population under 5. In the sensitivity analysis using IHME, early and late neonatal deaths (0--27 days) were combined for the $<$1 month band, the 1--5, 6--11 and 12--23 month groups were used directly, the 2--4 year group's deaths and person-time were divided equally across the three annual bands, and child-years were those implied by IHME's death counts and rates.

Because the attributable fraction of deaths in a mortality crisis (conflict, famine) need not resemble that of other years, crisis years took crisis-free values: a country-year whose all-cause under-five death rate in a source was more than 25\% above its crisis-free trend (the log-linear interpolation between the nearest earlier and later non-crisis years, found iteratively from the largest excess) took that trend in every age band. This applied to 14 country-years in four countries with UN IGME estimates (the Central African Republic, Malawi, Somalia and South Sudan) and two with IHME estimates (Somalia 2011 and South Sudan 2024). Averages across the 42 countries weight \textit{Pf}PR$_{2-10}$ by the UN WPP population under 5 and divide total deaths by total child-years. National totals are point estimates: because the seven age-band models were fitted separately, their sampling covariance is not estimated and age-specific confidence intervals were not summed. Each source's cause-specific malaria deaths were compared with the \textit{Pf}PR-ACM estimates based on the same source's all-cause mortality, over the same child-years (the UN WPP population for UN IGME; IHME person-years for IHME).


\subsection*{Sensitivity analyses}

We carried out sensitivity analyses to assess how the estimated relationship between prevalence and age-specific all-cause mortality depended on the subset of the data and on the handling of missing covariates and malnutrition. First, we refitted the seven age-band models in subsets of the primary sample: Sahelian survey regions (survey-region boundary centroid at or north of 12\textdegree N and west of 36\textdegree E, excluding the Horn of Africa; 31 surveys in 8 countries, 19,652 deaths), where transmission is highly seasonal; the UN M49 sub-regions of Western Africa (61 surveys in 15 countries, 56,426 deaths), Middle Africa (18 surveys in 7 countries, 17,628 deaths) and Eastern Africa (50 surveys in 12 countries, 33,671 deaths), with Southern Africa not refitted because \textit{Pf}PR$_{2-10}$ never exceeded about 10\%; surveys conducted up to (74 surveys) or after (62 surveys) the median survey year, 2014; and age bands entered up to or after the median calendar date of entry (late 2012). Subset fits used the same specification and covariate scaling, with knots at quantiles of the subset's own predictor values. Second, we imputed every remaining covariate gap rather than excluding records: survey-unit covariates by chained-equation multiple imputation (predictive mean matching, 10 imputations), and health expenditure (Somalia 2000--2012, South Sudan 2011--2016, Zimbabwe 2000--2009 and every country in 2024) and South Sudan's GDP after 2015 from generalised additive models of the observed national series. This retained 8,630,306 records (161 surveys in 38 countries, 122,092 deaths); imputation uncertainty was propagated by refitting with each of the ten imputed datasets and pooling with Rubin's rules. Third, we refitted the primary models without wasting and stunting (see Covariate adjustment).
